\documentclass[12pt]{article}
\usepackage[utf8]{inputenc}
\usepackage[T1]{fontenc}
\usepackage{hyperref}
\title{Taburta Language Resources for K-12 Curriculum Development}
\author{Compiled by Claude Sonnet 5 (claude-sonnet-5) for Andr\'as Kornai}
\date{}
\begin{document}
\maketitle

\section{Introduction}
Taburta (ISO 639-3 tbp, Glottocode tawo1244) is a Lakes Plain language of Papua, Indonesia. Search results show it under the name Taworta, with Dabra, Diebroud, Bok and Taria as other names. Speaker counts conflict across sources: about 300 (Wikipedia, 2006), 600 (Joshua Project, ASJP) and 150 (Global Recordings Network). Little documentation turned up in search.

\section{Primary Educational Resources}
No primers, textbooks or school materials were found in search results.

\section{Linguistic and Reference Materials}
\begin{itemize}
\item Wikipedia, Taworta language (overview, names, classification): \url{https://en.wikipedia.org/wiki/Taworta_language}
\item ASJP, Dabra/Taworta (word list entry, East Lakes Plain placement): \url{https://asjp.clld.org/languages/DABRA_TAWORTA}
\end{itemize}
The main source cited for the language is an unpublished 2006 SIL survey report by Rumaropen on the Diebrou language. I did not find that report online.

\section{Religious and Community Resources}
\begin{itemize}
\item Joshua Project, Taworta (Dabra) people group profile: \url{https://joshuaproject.net/people_groups/15303}
\item Joshua Project, language page for tbp: \url{https://joshuaproject.net/index.php/languages/tbp}
\end{itemize}

\section{Digital and Audio Resources}
\begin{itemize}
\item Global Recordings Network, language page for tbp (listing found in search; I did not check what recordings it holds): \url{https://globalrecordings.net/en/language/tbp}
\end{itemize}

\section{Curriculum Development Context}
The search found no evidence of schooling in Taburta, no orthography and no published teaching materials. A curriculum would have to start from basic documentation work with the community. The sources disagree on speaker numbers and location, so field verification is needed.
\end{document}